\section{Limitations}

We identify fifteen limitations of this study, organized by severity. Collectively, these limitations indicate that the reported results are directional, not definitive, and should be interpreted with appropriate caution.

\subsection{High Severity}

\begin{enumerate}
    \item \textbf{Small DDI test subgroup ($n=42$, 10 melanomas).} The Dark subgroup test set contains only 42 images, of which 10 are melanoma-positive. This small sample size yields wide bootstrap confidence intervals and limited statistical power. The paired Dark AUROC change from fine-tuning to augmentation was $-0.066$ (95\% bootstrap interval $[-0.151,\;0.008]$, $p=0.962$), so the controlled single-seed experiment provides no evidence of a fairness improvement. Subgroup conclusions should be treated as preliminary and require validation on larger cohorts.
\end{enumerate}

\subsection{Medium Severity}

\begin{enumerate}
    \setcounter{enumi}{1}
    \item \textbf{No external validation.} All experiments use a single dataset (DDI) for fine-tuning and evaluation. The generalizability of findings to other datasets, clinical settings, and populations is unknown. External validation on independent datasets is required before any clinical deployment claims.

    \item \textbf{Ablation tested saved-image quantity, not augmentation technique.} The ablation study varied the synthetic-image multiplier (0$\times$, 2$\times$, 5$\times$, and 10$\times$) rather than comparing different augmentation techniques. The non-monotonic pattern does not isolate the contribution of specific augmentation transforms.

    \item \textbf{Limited synthetic image quality validation.} A reproducible report now checks readability, output dimensions, pixel statistics, and edge-variance distributions (\texttt{results/synthetic\_quality.json}). These checks do not establish medical realism or lesion preservation; dermatologist review or a validated medical-image distribution metric such as FID is still required.

    \item \textbf{Sensitivity--specificity tradeoff across subgroups.} At the stage-specific validation-selected thresholds, the augmented model produced sensitivity of 0.600 in the Light subgroup versus 0.500 in the Dark test subgroup, but subgroup specificity and AUROC remained different. Threshold-dependent comparisons therefore do not establish equal clinical performance across groups.
\end{enumerate}

\subsection{Low Severity}

\begin{enumerate}
    \setcounter{enumi}{5}
    \item \textbf{Asymmetric oversampling with static pos\_weight.} We augmented only dark melanomas (29$\to$290) with zero dark benigns while pos\_weight stayed static at 2.9 (computed from the original training split). This incentivizes the model to associate dark skin tone with malignancy, which likely contributes to the Dark specificity collapse from 0.719 (Baseline) to 0.500 (Fine-tuned) to 0.469 (+Synthetic) and Dark AUROC below 0.5 (0.481), as discussed in Section 5.4. Future work should augment dark benign lesions symmetrically or recompute pos\_weight on the combined train+synthetic set.

    \item \textbf{Medium skin tone subgroup excluded.} The DDI dataset includes images across the Fitzpatrick scale, but the medium skin tone subgroup was excluded from per-subgroup analysis due to insufficient sample size. This exclusion limits the generalizability of findings to the full spectrum of skin tones.

    \item \textbf{MPS device non-determinism.} All experiments were conducted on Apple Silicon (MPS backend). The MPS backend introduces non-deterministic behavior in certain operations, which may contribute to minor variation across runs. Results should be replicated on CUDA-based hardware for confirmation.

    \item \textbf{Dynamic pos\_weight may differ across runs.} Pos\_weight is computed dynamically from the training split as $n_{\text{neg}}/n_{\text{pos}}$. While this is more principled than a hardcoded value, it means the exact pos\_weight depends on the split, which is seed-dependent. The split and pos\_weight are saved for reproducibility.

    \item \textbf{Single random seed.} The current reported run uses seed 42. While this ensures internal reproducibility, results from a single seed may not capture the full variance of the training process. Multi-seed evaluation is required for robust estimates.

    \item \textbf{Ablation figure Y-axis was initially mislabeled (corrected).} The ablation tradeoff figure was initially generated with incorrect axis labels during development. This error was identified and corrected. All figures in the manuscript have been regenerated and verified programmatically from the saved metrics (see Appendix).

    \item \textbf{Synthetic images use simple transforms.} The augmentation pipeline uses standard photometric and geometric transforms. More sophisticated approaches (e.g., diffusion models, GANs) may produce higher-quality synthetic images with better clinical realism, though at greater computational cost.

    \item \textbf{Bootstrap CI centers may not exactly match point estimates.} The bootstrap confidence interval centers do not exactly correspond to the reported point estimates, reflecting the asymmetric sampling distribution inherent in small subgroups. This is a known property of bootstrap methods and should be noted when interpreting CIs.

    \item \textbf{Synthetic images generated from only 29 source images.} With only 29 training dark-skin melanomas, the diversity of generated synthetics is limited by the source distribution. More source images would enable greater variation in the synthetic set.

    \item \textbf{Early stopping and single-seed training trajectory.} The controlled runs use a 10-epoch maximum with patience-three early stopping; the baseline, fine-tuning, and augmentation stages stopped before that maximum. This is part of the documented protocol, but the selected checkpoint remains sensitive to the validation split and random seed.
\end{enumerate}

\subsection{Fixed Issues}

The following issues from earlier versions have been addressed:

\begin{enumerate}
    \item \textbf{Synthetic images derived from training split only.} Synthetic augmentation images are now generated exclusively from training-split dark-skin melanoma images ($n\,{=}\,29$). No validation or test set images are used as sources. Byte-level SHA256 verification confirms zero overlap between synthetic and held-out images. This issue is resolved.
\end{enumerate}
